
\SGAProof{Proof}
If \(\dim H=1\), then clearly \(H=\mathbb G_a\).  Otherwise \(H\) has
dimension zero and is therefore an extension\footnote{Editor's note:
reference to Exposé VII\textsubscript{A}?} of an étale group \(H''\) by its
identity component \(H'\), which is radicial.  To describe \(H''_{\bar{k}}\),
it is enough to know the abstract group \(H''(\bar{k})\), which is isomorphic
to \(H(\bar{k})\).

The latter is a finite subgroup of \(\mathbb G_a(\bar{k})\).  It is therefore
zero when \(p=0\), and otherwise is of the form
\((\mathbb Z/p\mathbb Z)^r\), since it is annihilated by \(p\).  The group
\(H'\) is closed in \(\mathbb G_a\) and is defined by a single equation
(the coordinate ring \(\bar{k}[T]\) of \(\mathbb G_a\) is principal).
Since this equation has \(0\) as its only root over \(\bar{k}\), it is of the
form \(T^n=0\).  Compatibility with the group law implies that
\((T+T')^n\) belongs to the ideal generated by \(T^n\) and \(T'^n\) in
\(\bar{k}[T,T']\).  Thus, if \(p=0\), then \(n=1\) and \(H'\) is the trivial
group; while if \(p>0\), then \(n=p^m\) and
\(H'=\boldsymbol\alpha_{p^m}\).

The last assertion of \SGARefLink{sga3:xvii:proposition:1:5}{Proposition 1.5} follows more generally from the
following lemma.

\SGAStatement{Lemma 1.6}{%
If \(k\) is a perfect field, every extension \(H\) of an étale algebraic group
\(H''\) by a radicial group \(H'\) is trivial.  Moreover, \(H''\) has a unique
lifting in \(H\), namely \(H_{\mathrm{red}}\).}
\SGARefTarget{sga3:xvii:lemma:1:6}

Indeed, since \(k\) is perfect, the reduced \(k\)-scheme
\(H_{\mathrm{red}}\) is an algebraic subgroup of \(H\)
(Exposé VI\textsubscript{A}, \SGARefLink{sga3:VIA:0.2}{0.2}), is geometrically reduced and hence smooth
over \(k\) (Exposé VI\textsubscript{A}, \SGARefLink{prop:sga3-VIA-1.3.1}{1.3.1}), and is therefore étale because
\(H\) has dimension zero.  To see that the canonical projection
\(H_{\mathrm{red}}\to H''\) is an isomorphism, it is enough to check this after
the base-field extension \(k\to\bar{k}\); there it is enough to check that it
induces an isomorphism on \(\bar{k}\)-valued points, which is clear.  The final
assertion follows because every lifting of \(H''\) in \(H\), being étale over
\(k\), is reduced and hence necessarily contained in \(H_{\mathrm{red}}\).

Notice that \(\boldsymbol\alpha_{p^n}\) is an iterated extension of groups
isomorphic to \(\boldsymbol\alpha_p\).  \SGARefLink{sga3:xvii:proposition:1:5}{Proposition 1.5} therefore gives the
following corollary.

\SGAStatement{Corollary 1.7}{%
Let \(G\) be an algebraic group defined over an algebraically closed field
\(k\).  Then \(G\) is unipotent if and only if it has a composition series
whose successive quotients are isomorphic to \(\mathbb G_a\) when \(p=0\),
and to one of the groups
\(\mathbb G_a,\mathbb Z/p\mathbb Z,\boldsymbol\alpha_p\) when \(p>0\).
We shall call these the \emph{elementary unipotent groups}.}
\SGARefTarget{sga3:xvii:corollary:1:7}

\SGASection{2}{First properties of unipotent groups}
\SGARefTarget{sga3:xvii:section:2}

\SGAStatement{Proposition 2.1}{%
A unipotent algebraic group defined over a field \(k\) is affine over \(k\).}
\SGARefTarget{sga3:xvii:proposition:2:1}

By fpqc descent for affine morphisms, it suffices to prove the assertion when
\(k\) is algebraically closed.  In that case, \(G\) is by definition an
iterated extension of affine algebraic groups, and is therefore affine by
Exposé VI\textsubscript{B}, \SGARefLink{concl:sga3-VIB-9.2-viii}{9.2 (viii)}, applied to affine
morphisms.\footnote{Editor's note: verify this reference.}

\SGAStatement{Proposition 2.2}{%
\begin{enumerate}
\item[\rm\SGAInlineRefTarget{sga3:xvii:proposition:2:1:item:i}(i)] The property that an algebraic group is unipotent is invariant
  under extension of the base field.
\item[\rm\SGAInlineRefTarget{sga3:xvii:proposition:2:1:item:ii}(ii)] Every algebraic subgroup of a unipotent group is unipotent.
\item[\rm\SGAInlineRefTarget{sga3:xvii:proposition:2:1:item:iii}(iii)] Every algebraic quotient group of a unipotent group is
  unipotent.
\item[\rm\SGAInlineRefTarget{sga3:xvii:proposition:2:1:item:iv}(iv)] Every extension of a unipotent algebraic group by a unipotent
  algebraic group is likewise unipotent.
\end{enumerate}}
\SGARefTarget{sga3:xvii:proposition:2:2}

\SGAProof{Proof}
Assertion \SGARefLink{sga3:xvii:proposition:2:1:item:i}{(i)} follows immediately from \SGARefLink{sga3:xvii:definition:1:3}{Definitions 1.3} and \SGARefLink{sga3:xvii:proposition:1:2}{Proposition 1.2}.
To prove the remaining assertions, we may suppose that \(k\) is algebraically
closed.  Assertion \SGARefLink{sga3:xvii:proposition:2:1:item:iv}{(iv)} is then immediate from the definition.

Let \(G\) be a unipotent algebraic group, \(G'\) an algebraic subgroup of
\(G\), \(G''\) an algebraic quotient group of \(G\), and let
\(G_i\), \(1\leq i\leq n\), be a composition series of \(G\) such that
\[
  H_i=G_i/G_{i+1}
\]
is an elementary unipotent group (\SGARefLink{sga3:xvii:corollary:1:7}{Corollary 1.7}).

For \SGARefLink{sga3:xvii:proposition:2:1:item:ii}{(ii)}, consider the composition series of \(G'\) induced by that of \(G\):
\(G'_i=G_i\cap G'\).  The group \(G'_i/G'_{i+1}\) identifies with an
algebraic subgroup of \(H_i\), hence is isomorphic to a subgroup of
\(\mathbb G_a\).  It follows that \(G'\) is unipotent.

For \SGARefLink{sga3:xvii:proposition:2:1:item:iii}{(iii)}, consider the composition series of \(G''\) obtained as the image
of that of \(G\): let \(G''_i\) be the image of \(G_i\) in \(G''\).  Then
\(G''_i/G''_{i+1}\) is a quotient of \(H_i\), so it remains only to prove the
following lemma.

\SGAStatement{Lemma 2.3}{%
Let \(H\) be an elementary unipotent group (\SGARefLink{sga3:xvii:corollary:1:7}{Corollary 1.7}) defined over a
field \(k\).  Every algebraic quotient group \(H''\) of \(H\) is either zero
or isomorphic to \(H\) over \(k\).}
\SGARefTarget{sga3:xvii:lemma:2:3}

\SGAProof{Proof}
If \(H=\mathbb G_a\) in characteristic zero, or if
\(H=\boldsymbol\alpha_p\) or \(\mathbb Z/p\mathbb Z\) in characteristic
\(p>0\), it is enough to observe from \SGARefLink{sga3:xvii:proposition:1:5}{Proposition 1.5} that \(H\) has no
algebraic subgroups other than \(0\) and \(H\).  Indeed, a nonzero algebraic
subgroup of \(\boldsymbol\alpha_p\) is defined by a \(k\)-algebra of rank at
least \(p\) over \(k\), by \SGARefLink{sga3:xvii:proposition:1:5:item:iii}{Proposition 1.5(iii)}, and is therefore
\(\boldsymbol\alpha_p\).

Now let \(H=\mathbb G_a\) in characteristic \(p>0\), so that there is an exact
sequence
\[
  0\longrightarrow N\longrightarrow\mathbb G_a
    \longrightarrow H''\longrightarrow0.
\]
If \(N=\mathbb G_a\), then \(H''=0\).  Otherwise, arguing by induction on the
length of a composition series of \(N\), we may suppose that
\(N=\boldsymbol\alpha_p\), or that \(N\) is a form of
\((\mathbb Z/p\mathbb Z)^r\).

\begin{enumerate}
\item[\rm(a)] If \(N\simeq\boldsymbol\alpha_p\), the proof of
  Proposition 1.5(iii) shows that \(N\) is necessarily the kernel of the
  Frobenius morphism \(F\) of \(\mathbb G_a\), and the assertion follows from
  the exact sequence
  \[
    0\longrightarrow\boldsymbol\alpha_p\longrightarrow\mathbb G_a
      \xrightarrow{F}\mathbb G_a\longrightarrow0.
  \]
\item[\rm(b)] If \(N\simeq\mathbb Z/p\mathbb Z\), there is plainly an
  \(a\in k^\times\) such that \(N\) is the closed subscheme of
  \(\mathbb G_a=\operatorname{Spec}k[X]\) defined by
  \(X^p-aX=0\).  The assertion then follows from the exact sequence
  \[
    0\longrightarrow N\longrightarrow\mathbb G_a
      \xrightarrow{P}\mathbb G_a\longrightarrow0,
  \]
  where \(P\) is the Artin--Schreier morphism \(x\mapsto x^p-ax\).
\item[\rm(c)] If \(N\) is a form of \((\mathbb Z/p\mathbb Z)^r\), there is a
  finite Galois extension \(k'/k\) that trivializes \(N\).  By \SGARefLink{sga3:xvii:section:b}{(b)} and an
  evident induction on \(r\), the quotient \(\mathbb G_a/N\) is a form of
  \(\mathbb G_a\) trivialized by \(k'\).  It remains to apply the next lemma.
\end{enumerate}

\SGAStatement{Lemma 2.3 bis}{%
Let \(k\) be a field, and let \(G\) be a \(k\)-algebraic group that is a form
of \(\mathbb G_a\), trivialized by a finite separable extension \(k'/k\).
Then \(G\) is isomorphic to \(\mathbb G_a\).}
\SGARefTarget{sga3:xvii:lemma:2:3:bis}
